\documentclass[11pt]{article}
\usepackage[T1]{fontenc}
\usepackage{textcomp}
\usepackage[margin=0.75in]{geometry}
\usepackage{amsmath,amssymb,amsthm}
\usepackage{array,booktabs}
\usepackage{hyperref}
\setlength{\parindent}{0pt}
\newtheorem{theorem}{Theorem}
\theoremstyle{definition}
\newtheorem{definition}{Definition}
\title{Flow Proof Artifact: prop-01-equilateral-triangle}
\author{Generated by \texttt{flow doc proof}}
\date{Flow Proof Book}
\begin{document}
\maketitle
On a given finite straight line to construct an equilateral triangle.\\[0.5em]
\textbf{Source.} euclid — Elements, Book I, Proposition 1\\[0.5em]
\bigskip
\noindent{\large \textbf{Axiom 1.} \textit{On a given finite straight line to construct an equilateral triangle} \hfill the Euclidean plane \cdot Euclid Book I \cdot Proposition 1: equilateral triangle on a segment}\par
\smallskip\noindent\textit{Coordinate: the Euclidean plane \cdot Euclid Book I \cdot Proposition 1: equilateral triangle on a segment}\par
\smallskip\noindent\textit{Source: euclid — Elements, Book I, Proposition 1}\par
\medskip
\noindent\textbf{Goal.} On a given finite straight line to construct an equilateral triangle\par
\noindent$\displaystyle \text{equilateral on AB exists}$\par
\medskip
\noindent
\begin{tabular}{@{} >{\raggedright\arraybackslash}p{0.06\textwidth} >{\raggedright\arraybackslash}p{0.47\textwidth} >{\raggedleft\arraybackslash}p{0.41\textwidth} @{}}
\textbf{\#} & \textbf{Proof} & \textbf{Mathematics} \\
\midrule
\textbf{1.} & We accept proposition 1: equilateral triangle on a segment for Book I of the Elements in the Euclidean plane without proof — an ontological commitment, not a lemma. & \\[0.45em]
\textbf{2.} & Let C = intersection\_of\_circles(center\_A, radius\_AB, center\_B, radius\_BA). & \\[0.45em]
\textbf{3.} & We invoke the derived fact: Postulate 3: a circle can be drawn with any center and radius. & \\[0.45em]
\textbf{4.} & We invoke the derived fact: Postulate 1: a straight line can be drawn from any point to any point. & \\[0.45em]
\textbf{5.} & From step 2, step 3, and step 4, we can deduce that segment AC equals segment AB. & $\displaystyle \text{segment AC} = \text{segment AB}$ \\[0.45em]
\textbf{6.} & This holds immediately by the stated axiom: segment BC equals segment AB. & $\displaystyle \text{segment BC} = \text{segment AB}$ \\[0.45em]
\textbf{7.} & This holds immediately by the stated axiom: equilateral on AB exists. Hence proven. & $\displaystyle \text{equilateral on AB exists}$ \\[0.45em]
\bottomrule
\end{tabular}
\label{thm:Geometry-Euclid Book I-proposition_1_equilateral_triangle_on_a_segment}
\par\medskip\hrule
\end{document}
